% Insertion fragment; no preamble or document environment.
% Mathematical dependency: Theorem \ref{sgd:parametrization}, whose label
% was verified in agent_work/two_cross_scalar_gauss_domain_20260905/
% scalar_gauss_domain.tex, lines 214--215, on 5 September 2026.
% Insert after that theorem's complete proof and before the section
% "A short exact family with every equation checked" (currently line 341).
% The fragment uses only the parent article's amsmath/amssymb/amsthm
% packages and its existing theorem and lemma environments.
% All new labels have the prefix sgd:literal-addon:.

\section{The complete literal scalar--Gauss domain and its phase quotient}
\label{sgd:literal-addon:section}

The classified-frame bijection of Theorem~\ref{sgd:parametrization}
has an exact extension to the entire literal two-cross domain.
We give the forced temporal classification, the full parameter map,
and the complete fibres of that map. All spaces retain their underlying
real structures, $[X,Y]=XY-YX$, and $*$ denotes the adjoint.

\subsection{The literal domain and the retained constants}

Retain exactly
\[
 H_0=\begin{pmatrix}0&0&0\\0&-1/2&1/(2\sqrt2)\\
 0&1/(2\sqrt2)&1/2\end{pmatrix},\qquad
 q=\frac1{2\sqrt2}\begin{pmatrix}1\\0\\0\end{pmatrix},\qquad
 A_0=\sqrt2\begin{pmatrix}0&-1&0\\1&0&0\\0&0&0\end{pmatrix},
\]
\[
 m^2=\Gamma=1,\qquad \kappa=6,\qquad
 F_0=[A_0,[A_0,H_0]]=-2H_0+\operatorname{diag}(-2,1,1).
\]
Here $\mathfrak{su}(3)$ consists of the traceless anti-Hermitian complex
$3\times3$ matrices, with positive real inner product
$\langle X,Y\rangle=-\operatorname{tr}(XY)$.
For $Y_{\rm tf}=Y-(\operatorname{tr}Y)I_3/3$, define
\[
 \mathcal I_{H_0}Y
 =([H_0,[H_0,Y]]+qq^*Y+Yqq^*)_{\rm tf},\qquad
 \mathcal K_{Z_1,Z_2}Y=-\sum_{j=1}^2[Z_j,[Z_j,Y]].
\]
Multiplication of the retained matrices gives
$H_0q=0$, $H_0^2=(3/8)\operatorname{diag}(0,1,1)$ and $q^*q=1/8$.
Consequently $\operatorname{tr}H_0^2+2q^*q=1$, so the original
scalar coefficient $m^2+\Gamma(\operatorname{tr}H_0^2+2q^*q)$ is $2$.
Also $\operatorname{adj}H_0=\operatorname{diag}(-3/8,0,0)$ gives
$6(\operatorname{adj}H_0-qq^*)_{\rm tf}
=\operatorname{diag}(-2,1,1)$, agreeing with the displayed force.

Put
\[
 P=\frac{qq^*}{q^*q}=\operatorname{diag}(1,0,0),\qquad
 \mathfrak p=\{X\in\mathfrak{su}(3):
 PXP=(I_3-P)X(I_3-P)=0\}.
\]
Let $\mathcal S_{\rm sg}$ be exactly the set of tuples
\[
 (t,\omega,C_0,Z_{1,0},Z_{2,0}),\qquad
 t=g^2>0,\quad \omega\in\mathbb R\setminus\{0\},\quad
 C_0\in\mathfrak{su}(3),\quad Z_{j,0}\in\mathfrak p,
\]
for which the following three full residuals vanish:
\begin{align}
 [C_0,[C_0,H_0]]
 -\sum_j[Z_{j,0},[Z_{j,0},H_0]]-F_0&=0,
 \label{sgd:literal-addon:H}\\
 (C_0^2-Z_{1,0}^2-Z_{2,0}^2)q+2q&=0,
 \label{sgd:literal-addon:q}\\
 (t\mathcal I_{H_0}+2\mathcal K_{Z_{1,0},Z_{2,0}})C_0
 -4\omega[Z_{1,0},Z_{2,0}]&=0.
 \label{sgd:literal-addon:G}
\end{align}
Their codomains are, respectively,
$\operatorname{Herm}_3(\mathbb C)_0$, $\mathbb C^3$, and
$\mathfrak{su}(3)$. The original Gauss residual is
\[
 (\mathcal I_{H_0}+2\mathcal K_{Z_{1,0},Z_{2,0}}/g^2)C_0
 -(4\omega/g^2)[Z_{1,0},Z_{2,0}].
\]
Multiplication by the same retained positive $t=g^2$ gives
\eqref{sgd:literal-addon:G}, and division by $t$ is its inverse.
No spatial residual is required to vanish in the definition of
$\mathcal S_{\rm sg}$.

\subsection{Temporal classification and all inverse frame maps}

For $v\in\mathbb C^2$, define
\[
 X(v)=\begin{pmatrix}0&-v^*\\v&0_2\end{pmatrix},\qquad
 \mathfrak k=\{X\in\mathfrak{su}(3):[X,P]=0\}.
\]
Every element of $\mathfrak p$ has exactly this form, whose inverse
is the lower first column. Define the fixed constants and matrix
\[
 r=\frac{\sqrt6}{4},\qquad \delta=\frac1{2r}=\sqrt{\frac23},
 \qquad x=\frac12,\qquad d_0=\frac1{2\sqrt2},
\]
\[
 S=\frac1{\sqrt{2r(r+x)}}
 \begin{pmatrix}d_0&-(r+x)\\r+x&d_0\end{pmatrix},
 \qquad U=\operatorname{diag}(1,S).
\]
The identity $r^2=x^2+d_0^2$ shows that the numerator columns
are orthogonal, have squared length $2r(r+x)$, and have determinant
$2r(r+x)>0$. Multiplication by the lower block of $H_0$ gives
their respective eigenvalues $r,-r$. Thus
\[
 U^*U=I_3,\qquad \det U=1,\qquad
 U^*H_0U=\operatorname{diag}(0,r,-r),\qquad Uq=q.
\]

\begin{lemma}[The temporal form is forced by scalars and Gauss]
\label{sgd:literal-addon:temporal}
Every tuple in $\mathcal S_{\rm sg}$ has $C_0\in\mathfrak k$.
In the fixed first frame $C_d=U^*C_0U$, there are unique
$e\in\mathbb R\setminus\{0\}$ and $z\in\mathbb C$ with
$|z|=|e|/2$ such that
\[
 C_d=\begin{pmatrix}
 ie&0&0\\0&-ie/2&z\\0&-\bar z&-ie/2
 \end{pmatrix}.
\]
Set $\gamma=-2iz/e$. Then $|\gamma|=1$. For either square root
$\zeta^2=-\bar\gamma$, define
\[
 W(\zeta)=\frac{\zeta}{\sqrt2}
 \begin{pmatrix}\gamma&\gamma\\1&-1\end{pmatrix}
 =\frac1{\sqrt2}
 \begin{pmatrix}-\zeta^{-1}&-\zeta^{-1}\\\zeta&-\zeta\end{pmatrix},
 \qquad B(\zeta)=U\operatorname{diag}(1,W(\zeta)).
\]
The simultaneous change and its inverse are
\begin{equation}
 Y'=B(\zeta)^*Y_0B(\zeta),\qquad
 Y_0=B(\zeta)Y'B(\zeta)^*.
 \label{sgd:literal-addon:frame}
\end{equation}
They give
\[
 C'=ie\operatorname{diag}(1,0,-1),\qquad
 H'=\begin{pmatrix}0&0&0\\0&0&r\\0&r&0\end{pmatrix},\qquad
 q'=q,\qquad F'=-2H'+\operatorname{diag}(-2,1,1).
\]
These maps preserve the spatial class and intertwine all three full
residuals in \eqref{sgd:literal-addon:H}--\eqref{sgd:literal-addon:G}.
\end{lemma}

\begin{proof}
Block multiplication gives the orthogonal direct sum
$\mathfrak{su}(3)=\mathfrak k\oplus\mathfrak p$, together with
$[\mathfrak p,\mathfrak p]\subset\mathfrak k$ and
$[\mathfrak k,\mathfrak p]\subset\mathfrak p$.
Since $H_0,qq^*$ are block diagonal, $\mathcal I_{H_0}$ preserves
both summands. More explicitly, $H_0X(v)H_0=0$; the two $H_0^2$
products sum to $(3/8)X(v)$, and the two $qq^*$ products sum
to $X(v)/8$. Each term is traceless. Thus
\[
 \mathcal I_{H_0}|_{\mathfrak p}
 =\tfrac12\operatorname{id}_{\mathfrak p}.
\]
The two bracket inclusions make $\mathcal K=\mathcal K_{Z_{1,0},Z_{2,0}}$
preserve both summands as well. Cyclicity gives, for every
$Y\in\mathfrak{su}(3)$,
\[
 \langle Y,\mathcal KY\rangle
 =-\sum_j\operatorname{tr}([Z_{j,0},Y]^2)
 =\sum_j\|[Z_{j,0},Y]\|^2\ge0.
\]
The $\mathfrak p$ projection of Gauss is therefore
\[
 (\tfrac t2\operatorname{id}_{\mathfrak p}
       +2\mathcal K|_{\mathfrak p})(C_0)_{\mathfrak p}=0.
\]
Pairing with $(C_0)_{\mathfrak p}$ gives nonnegative summands,
including $(t/2)\|(C_0)_{\mathfrak p}\|^2$. Since $t>0$,
$(C_0)_{\mathfrak p}=0$, proving $C_0\in\mathfrak k$.

The arbitrary first-frame temporal block consequently has the form
\[
 C_d=\begin{pmatrix}
 ia&0&0\\0&i(h-a)/2&z\\0&-\bar z&-i(a+h)/2
 \end{pmatrix},\qquad a,h\in\mathbb R,\quad z\in\mathbb C.
\]
Write the spatial matrices in this first frame as
$X(\widetilde v),X(\widetilde w)$, and define
$\widetilde M=\widetilde v\widetilde v^*
                +\widetilde w\widetilde w^*$.
The whole vector residual is
\[
 (-a^2+\operatorname{tr}\widetilde M+2)q,
\]
and the whole adjoint residual is
\[
 \begin{pmatrix}
 2-2r(\widetilde M_{11}-\widetilde M_{22})&0&0\\
 0&-4r|z|^2+2r\widetilde M_{11}-1+2r&-2irhz\\
 0&2irh\bar z&4r|z|^2-2r\widetilde M_{22}-1-2r
 \end{pmatrix}.
\]
For completeness, if $D=\operatorname{diag}(r,-r)$, the lower
spatial contribution is $\widetilde MD+D\widetilde M$, whose
diagonal entries are $2r\widetilde M_{11},-2r\widetilde M_{22}$
and whose off-diagonal entries vanish. The temporal double
commutator has diagonal entries $-4r|z|^2,4r|z|^2$ and
off-diagonal entries $-2irhz,2irh\bar z$.
Subtracting the lower force $I_2-2D$ gives the displayed lower block.
The upper spatial contribution and force give
$2-2r(\widetilde M_{11}-\widetilde M_{22})$.
Finally $X(v)^2=\operatorname{diag}(-v^*v,-vv^*)$ gives the
displayed vector residual, with all its other components zero.

Thus both scalar equations are equivalent to
\[
 a^2=2+\operatorname{tr}\widetilde M,\qquad
 2|z|^2-\widetilde M_{11}=1-\delta,\qquad
 2|z|^2-\widetilde M_{22}=1+\delta,\qquad hz=0.
\]
The upper adjoint entry is minus the lower trace.
Adding the middle equations gives
$4|z|^2=2+\operatorname{tr}\widetilde M=a^2$.
Since $\widetilde M_{22}\ge0$ and $1+\delta>0$, $z\ne0$.
Hence $h=0$ and $a\ne0$. This proves the claimed temporal block
with $e=a$, where its entries recover $e,z$ uniquely.

Now $|\gamma|=1$ and each root $\zeta^2=-\bar\gamma$ has
$|\zeta|=1$, with $\gamma=-\zeta^{-2}$. The displayed entries give
$W^*W=I_2$ and $\det W=-\zeta^2\gamma=1$.
The lower Hermitian block of $-iC_d$ is
\[
 \frac e2\begin{pmatrix}-1&\gamma\\\bar\gamma&-1\end{pmatrix}.
\]
Its products with $(\gamma,1)^{\mathsf T}$ and
$(\gamma,-1)^{\mathsf T}$ have eigenvalues $0,-e$,
respectively. The same columns give
$W^*DW=\left(\begin{smallmatrix}0&r\\r&0\end{smallmatrix}\right)$.
This proves the stated classified matrices.

The block-diagonal unitary $B(\zeta)$ fixes $q$ and commutes with
$P$. Its conjugation therefore preserves $\mathfrak p$ and
$\mathfrak k$. Unitarity gives
$B^*(XY)B=(B^*XB)(B^*YB)$; adjoints, commutators and trace
projection satisfy the corresponding equalities.
It follows that $\mathcal I,\mathcal K$ and every full matrix
residual are intertwined by conjugation. The vector residual is
multiplied by $B^*$. Equation~\eqref{sgd:literal-addon:frame} gives
the exact inverse on every datum, with no rescaling.
\end{proof}

\subsection{The full parameter map and the quotient bijection}

We restate the classified parameters so that the literal map is explicit.
Put
\[
 Q=\operatorname{diag}(1,2),\qquad E=\operatorname{diag}(2/3,-5/6),
\]
\[
 \mathcal D=\{(u,v_0,y)\in\mathbb R^3:
 u>0,\ d=uv_0-\delta^2-y^2>0,\ c=4v_0-5u>0\}.
\]
For $(u,v_0,y)\in\mathcal D$, independent
$\epsilon,\sigma\in\{1,-1\}$, and $p_+,p_-\in U(1)$, define
\begin{align}
 m&=\delta+iy,&
 M&=\begin{pmatrix}u&m\\\bar m&v_0\end{pmatrix},&
 t&=\frac{72d}{c},\notag\\
 T&=\frac t4E-\frac12(QM+MQ),&
 k&=\sigma\sqrt{\frac d{-\det T}},&
 e&=\epsilon\sqrt{2+u+v_0},\notag\\
 \omega&=e/k,&N&=kT,&A_\pm&=M\pm N.
 \label{sgd:literal-addon:parameters}
\end{align}
Theorem~\ref{sgd:parametrization} proves that $\det T<0$,
that $A_\pm$ are positive semidefinite of rank one, and that
$h_\pm=(A_\pm)_{11}>0$. Thus the following vectors are defined
on the whole displayed domain:
\begin{equation}
 a_\pm=\begin{pmatrix}
 \sqrt{h_\pm}\\(A_\pm)_{21}/\sqrt{h_\pm}
 \end{pmatrix},\qquad
 v=\frac{p_+a_++p_-a_-}{2},\qquad
 w=\frac{i(p_+a_+-p_-a_-)}{2}.
 \label{sgd:literal-addon:vectors}
\end{equation}
In particular $vv^*+ww^*=M$ and $i(vw^*-wv^*)=N$.
The cited theorem proves that these parameters are in bijection
with all solutions of the two full scalar equations and full Gauss
in the frame $C'=ie\operatorname{diag}(1,0,-1)$,
$H'_{23}=H'_{32}=r$, $q'=q$, and
$F'=-2H'+\operatorname{diag}(-2,1,1)$.
It also proves the exact inverse phases
\begin{equation}
 p_+=\frac{(v-iw)_1}{\sqrt{h_+}},\qquad
 p_-=\frac{(v+iw)_1}{\sqrt{h_-}}.
 \label{sgd:literal-addon:phase-inverse}
\end{equation}

Define the full raw parameter space
\[
 \mathcal P_{\rm lit}
 =\mathcal D\times\{1,-1\}^2\times U(1)^2\times U(1)
\]
with coordinates $(u,v_0,y,\epsilon,\sigma,p_+,p_-,\zeta)$.
For the last unit phase set $\gamma=-\zeta^{-2}$ and use the
matrix $B(\zeta)$ of Lemma~\ref{sgd:literal-addon:temporal}.
Define the map $\mathfrak F:\mathcal P_{\rm lit}\to\mathcal S_{\rm sg}$
by \eqref{sgd:literal-addon:parameters}--\eqref{sgd:literal-addon:vectors}
and the following entire literal tuple:
\begin{equation}
 \begin{gathered}
 \mathfrak F(u,v_0,y,\epsilon,\sigma,p_+,p_-,\zeta)
 =(t,\omega,C_0,Z_{1,0},Z_{2,0}),\\
 \begin{aligned}
 C_0&=B(\zeta)\bigl(ie\operatorname{diag}(1,0,-1)\bigr)B(\zeta)^*,\\
 Z_{1,0}&=B(\zeta)X(v)B(\zeta)^*,\\
 Z_{2,0}&=B(\zeta)X(w)B(\zeta)^*.
 \end{aligned}
 \end{gathered}
 \label{sgd:literal-addon:forward}
\end{equation}
The codomain assertion follows from the classified theorem and the
exact residual intertwining in Lemma~\ref{sgd:literal-addon:temporal}.

\begin{theorem}[Exact quotient parameterization of the entire literal domain]
\label{sgd:literal-addon:quotient}
The map $\mathfrak F$ is onto $\mathcal S_{\rm sg}$.
The map
\begin{equation}
 \iota(u,v_0,y,\epsilon,\sigma,p_+,p_-,\zeta)
 =(u,v_0,y,\epsilon,\sigma,-p_+,-p_-,-\zeta)
 \label{sgd:literal-addon:involution}
\end{equation}
is a free involution of $\mathcal P_{\rm lit}$.
Each fibre of $\mathfrak F$ is exactly one two-element
$\iota$-orbit. Consequently \eqref{sgd:literal-addon:forward}
induces a bijection of sets
\begin{equation}
 \overline{\mathfrak F}:
 \mathcal P_{\rm lit}/\langle\iota\rangle
 \ \longrightarrow\ \mathcal S_{\rm sg}.
 \label{sgd:literal-addon:bijection}
\end{equation}
The inverse of this bijection is recovered from the literal matrices
by the formulas in the proof, with both square-root representatives
retained.
\end{theorem}

\begin{proof}
We first construct a preimage of every literal tuple, recording all
inverse coordinates. Given
$(t,\omega,C_0,Z_{1,0},Z_{2,0})\in\mathcal S_{\rm sg}$, compute
\[
 C_d=U^*C_0U,\qquad
 e=-i(C_d)_{11},\qquad z=(C_d)_{23},\qquad
 \gamma=-2iz/e.
\]
Lemma~\ref{sgd:literal-addon:temporal} proves $e\in\mathbb R\setminus\{0\}$
and $|\gamma|=1$. Choose either square root
$\zeta^2=-\bar\gamma$. Define
\[
 Z'_j=B(\zeta)^*Z_{j,0}B(\zeta),\qquad
 v=\begin{pmatrix}(Z'_1)_{21}\\(Z'_1)_{31}\end{pmatrix},\qquad
 w=\begin{pmatrix}(Z'_2)_{21}\\(Z'_2)_{31}\end{pmatrix}.
\]
The lemma proves that these are classified solutions with the original
$t,\omega,e$. Compute
\[
 M=vv^*+ww^*,\qquad N=i(vw^*-wv^*),\qquad
 u=M_{11},\quad v_0=M_{22},\quad y=\operatorname{Im}M_{12},
\]
\[
 \epsilon=\operatorname{sign}e,\qquad
 \sigma=\operatorname{sign}(e/\omega).
\]
Theorem~\ref{sgd:parametrization} proves
$(u,v_0,y)\in\mathcal D$, $M_{12}=\delta+iy$, and all equations
in \eqref{sgd:literal-addon:parameters} with this unchanged
$t,\omega,e$. It proves $h_\pm=(M\pm N)_{11}>0$, so
\eqref{sgd:literal-addon:phase-inverse} recovers unique unit
$p_+,p_-$. The vector formulas recover the displayed $v,w$
and the inverse frame map then recovers all three literal matrices.
Thus \eqref{sgd:literal-addon:forward} sends the recovered raw
parameter to the original tuple. This proves surjectivity.

Next we prove that \eqref{sgd:literal-addon:involution} preserves
each fibre. It preserves the parameter domain, and applying it
twice is the identity. It has no fixed point because
$\zeta=-\zeta$ would imply $\zeta=0$, contrary to $|\zeta|=1$.
Replacing $\zeta$ by $-\zeta$ preserves $\gamma=-\zeta^{-2}$ and
gives $W(-\zeta)=-W(\zeta)$. Put
$D_0=\operatorname{diag}(1,-I_2)$. Then
 \begin{align*}
 B(-\zeta)&=B(\zeta)D_0,\\
 D_0\bigl(ie\operatorname{diag}(1,0,-1)\bigr)D_0
 &=ie\operatorname{diag}(1,0,-1),\\
 D_0X(v)D_0&=X(-v).
 \end{align*}
Replacing $(p_+,p_-)$ by $(-p_+,-p_-)$ replaces $(v,w)$ by
$(-v,-w)$ in \eqref{sgd:literal-addon:vectors}, while leaving
$M,N,t,e,k,\omega$ fixed. The preceding matrix equalities
therefore give $\mathfrak F\circ\iota=\mathfrak F$ on the entire
raw parameter domain. This proves that $\overline{\mathfrak F}$
is well defined.

To determine every fibre, suppose two raw parameters have the same
literal image. The common temporal matrix $C_d$ recovers the same
$e,z,\gamma$ by the inverse formulas above. Their frame roots
therefore satisfy either $\zeta_2=\zeta_1$ or
$\zeta_2=-\zeta_1$, since they square to the same nonzero complex
number. If $\zeta_2=\zeta_1$, the classified temporal and spatial
matrices obtained by that common inverse frame map coincide.
The classified bijection of Theorem~\ref{sgd:parametrization}
then forces equality of all other raw parameters.

If $\zeta_2=-\zeta_1$, the identity
$B(-\zeta)=B(\zeta)D_0$ implies
\[
 Z'_{j,2}=D_0Z'_{j,1}D_0,\qquad
 v_2=-v_1,\qquad w_2=-w_1.
\]
Hence $M_2=M_1$ and $N_2=N_1$, so the inverse coordinates
$u,v_0,y$ are equal. The original $e,\omega$ are equal, which
also gives equal $\epsilon,\sigma$. Consequently $h_\pm$ agree,
and \eqref{sgd:literal-addon:phase-inverse} gives
\[
 p_{+,2}=-p_{+,1},\qquad p_{-,2}=-p_{-,1}.
\]
Thus the two raw parameters differ by exactly $\iota$.
No further case is possible. Each fibre is therefore precisely
one two-element orbit. This proves injectivity of
$\overline{\mathfrak F}$ as well as its already-proved surjectivity.

Finally the inverse formulas starting from a literal tuple give
the same orbit for either root choice, by the second fibre case.
Starting from a raw parameter, applying $\mathfrak F$ and then those
inverse formulas gives its orbit, by the two fibre cases.
Starting from a literal tuple, applying the inverse formulas and
then $\overline{\mathfrak F}$ gives that literal tuple by the first
paragraph. These prove both compositions are identities on their
stated domains, completing the quotient bijection.
\end{proof}

The construction retains
$\operatorname{sign}e=\epsilon$ and
$\operatorname{sign}\omega=\epsilon\sigma$.
The involution changes neither sign. The two independent vector phases
and the frame phase remain coordinates, subject only to the proved
simultaneous-sign identification.

For explicit scope, the additional full spatial residuals would be
\begin{align*}
 R_1={}&-\omega^2Z_{1,0}-2\omega[C_0,Z_{2,0}]
 +[C_0,[C_0,Z_{1,0}]]
 +\mathcal K_{Z_{1,0},Z_{2,0}}Z_{1,0}
 +(t/2)\mathcal I_{H_0}Z_{1,0},\\
 R_2={}&-\omega^2Z_{2,0}+2\omega[C_0,Z_{1,0}]
 +[C_0,[C_0,Z_{2,0}]]
 +\mathcal K_{Z_{1,0},Z_{2,0}}Z_{2,0}
 +(t/2)\mathcal I_{H_0}Z_{2,0}.
\end{align*}
Neither $R_1=0$ nor $R_2=0$ is part of
$\mathcal S_{\rm sg}$ or used in either proof above.
The bijection therefore classifies exactly the two scalar equations
and Gauss; it does not assert a solution of these two additional
spatial equations.
